This section describes the experimental methodology: how we built the evaluation dataset, how we run the benchmark, and how we measure performance.

\subsection{Benchmark Dataset Construction}

A core challenge in this work was constructing a test set that actually requires retrieval to answer correctly. If we created generic compliance questions about widely-known regulations, a language model could answer them from its training data alone — making it impossible to measure the contribution of the retrieval component. We therefore designed a \textit{RAG-aware} test set.

To build the dataset, we identified specific passages in the YÖK regulatory documents that contain concrete details: exact article numbers, numerical thresholds (e.g., minimum GPA requirements, maximum study durations), and procedural rules. For each selected passage, we generated a synthetic university procedure document that either correctly follows the regulation (compliant), partially follows it (partial), or clearly violates a specific numerical or procedural requirement (non-compliant). The violation in non-compliant cases was always a concrete change — for example, changing a required minimum score from 55 to 45, or shortening a mandated notice period.

Gold labels were assigned by running GPT-4o \textit{with the original YÖK passage as context} (i.e., oracle evaluation). This ensures that the gold label reflects the actual regulatory requirement rather than the model's general intuition. Cases where the oracle and the generated label disagreed were discarded.

The final dataset contains 40 cases: 13 compliant, 14 partial, and 13 non-compliant. Cases span nine source documents and cover diverse topic areas including graduate admission, GPA requirements, financial aid, international student rules, and minor program enrollment.

\subsection{Evaluation Pipelines}

We evaluate five retrieval pipeline configurations, each combined with two LLM backends (GPT-4o and GPT-4o-mini), for a total of ten system configurations. Each configuration is evaluated on all 40 test cases, giving 400 LLM calls in total for the main benchmark.

All pipelines use the same document collection and chunking parameters (350-word chunks, 59-word overlap, 248 total chunks). For retrieval pipelines, the default retrieval depth is Top-$K = 5$, which was identified as optimal in our ablation study (Section~\ref{sec:results}).

\subsection{Evaluation Metrics}

\paragraph{Classification Accuracy.} Our primary metric is exact-match accuracy: the fraction of test cases where the system's predicted label (compliant / partial / non-compliant) matches the gold label.

\paragraph{Per-Class Accuracy.} Because the task is a three-class classification problem, we also report per-class accuracy (recall per class) to detect label bias. This is important because, as we discuss in Section~\ref{sec:results}, all models show a strong preference for the \textit{partial} label.

\paragraph{Latency.} We measure mean wall-clock latency per request (seconds), including the API round-trip time.

\paragraph{Cost.} We compute the total API cost for evaluating all 40 cases per configuration, using OpenAI's published pricing (GPT-4o: \$2.50/1M input tokens, \$10.00/1M output tokens; GPT-4o-mini: \$0.15/1M input tokens, \$0.60/1M output tokens at time of writing).

\subsection{Ablation Studies}

We conduct two ablation studies using the BM25 pipeline with GPT-4o-mini, which is a cost-effective configuration that allows us to run many experiments without large API costs.

\paragraph{Chunk Size Ablation.} We vary the chunk size across four values: 150, 250, 350, and 500 words, while keeping all other parameters fixed (Top-$K = 5$, overlap = 17\% of chunk size). For each chunk size, we rebuild the BM25 index from scratch and evaluate on the full 40-case test set.

\paragraph{Top-K Ablation.} We fix the chunk size at 350 words and vary the retrieval depth $K \in \{1, 3, 5, 10\}$. This measures how sensitive performance is to the number of retrieved passages provided as context to the LLM.
